\documentclass{article}
\usepackage{../math_template}

\author{Jonathan Kasongo}
\title{Miscellaneous Theorems and Ideas}

\begin{document}
\maketitle
\textit{This is just a collection of some cool facts and theorems I have
found, along with their proofs.}


\begin{problem}{Theorem 1.}
For an odd positive integer $n$, and positive integers $a,b$ with
$a^2 \leq b < (a+1)^2$ we have the following,

$$
\floor{\left( a + \sqrt{b} \right)^n} =
\left( a + \sqrt{b} \right)^n + \left( a - \sqrt{b} \right)^n.
$$
\end{problem}

\begin{solution}{Proof of Theorem 1.}
Firstly, due to Binomial expansion,

\[
\begin{aligned}
\left( a + \sqrt{b} \right)^n + \left( a - \sqrt{b} \right)^n &=
\sum_{r=0}^n \binom{n}{r} \left(\sqrt{b}\right)^r a^{n-r} +
\sum_{r=0}^n \binom{n}{r} \left(-\sqrt{b}\right)^r a^{n-r} \\
&= \sum_{r=0}^n \binom{n}{r} \left( \sqrt{b} \right)^r \left[ 1 + (-1)^r \right] a^{n - r}\\
&= 2 \sum_{\substack{r=0 \\ r \ \text{even}}}^{n} \binom{n}{r} \left(\sqrt{b}\right)^r
a^{n-r}
\end{aligned}
\]

which is clearly an integer. Now it suffices to show that

\[
\left(a + \sqrt{b}\right)^n - 1 <
\left(a + \sqrt{b}\right)^n + \left(a - \sqrt{b}\right)^n
\leq \left(a + \sqrt{b}\right)^n
\]

That is

$$
-1 < \left(a - \sqrt{b}\right)^n \leq 0
$$

but we know that $a \leq \sqrt{b} < a + 1$ and so
$a - \sqrt{b} \leq 0$ and $-1 < a - \sqrt{b}$ and raising both inequalites
to the power of $n$ gives the desired result. $\blacksquare$
\end{solution}
\newpage


\begin{problem}{Theorem 2.}
Let $C_1$ and $C_2$ be externally tangent circles with centres $O_1, O_2$
and radii $r_1, r_2, (r_1 \geq r_2)$ respectively. Let $P_1, P_2$ be points on $C_1, C_2$
respectively such that $P_1 P_2$ is a common tangent to both circles.
\vspace{0.2cm}

(a) Show that $P_1P_2 = 2\sqrt{r_1 r_2}$.\vspace{0.2cm}

(b) Find the area of the region bounded by the two circles and their common
tangent.
\end{problem}

\begin{solution}{Solution}
(a) The idea is to draw in the line $l$ through $P_2$ that is parallel to
$O_1 O_2$. Let $Q = l \cap O_1P_1$.

\begin{center}
\begin{asy}
/* Geogebra to Asymptote conversion, documentation at artofproblemsolving.com/Wiki go to User:Azjps/geogebra */
import graph; size(250);
real labelscalefactor = 0.5; /* changes label-to-point distance */
pen dps = linewidth(0.7) + fontsize(10); defaultpen(dps); /* default pen style */
pen dotstyle = black; /* point style */
real xmin = -4.352660005815247, xmax = 9.312262465894712, ymin = -5.495146144789845, ymax = 7.539749049794846;  /* image dimensions */
pen qqzzff = rgb(0.,0.6,1.); pen ffwwqq = rgb(1.,0.4,0.);
pair O_1 = (0.,0.), O_2 = (6.,0.), P_1 = (1.3333333333333333,3.7712361663282534), P_2 = (6.666666666666667,1.8856180831641276), Q = (0.666666666666667,1.8856180831641276);

filldraw(circle(O_1, 4.), qqzzff + opacity(0.05000000074505806),  qqzzff);
filldraw(circle(O_2, 2.), qqzzff + opacity(0.05000000074505806),  qqzzff);
filldraw((1.2304643290459811,3.4802786843062745)--(1.52142181106796,3.377409680018922)--(1.6242908153553124,3.668367162040901)--P_1--cycle, ffwwqq + opacity(0.10000000149011612),  ffwwqq);
/* draw figures */
draw((xmin, -0.35355339059327373*xmin + 4.242640687119285)--(xmax, -0.35355339059327373*xmax + 4.242640687119285),  ffwwqq); /* line */
draw(P_1--O_1,  ffwwqq);
draw(P_2--O_2,  ffwwqq);
draw(Q--P_2,  qqzzff);
draw(O_1--(4.,0.),  qqzzff);
draw(O_2--(4.,0.),  qqzzff);
label("$r_1$",(2.041463350518847,-0.13488113233546092),SE*labelscalefactor);
label("$r_2$",(5.0383254153452315,-0.13488113233546092),SE*labelscalefactor);
label("$r_1 + r_2$",(2.992580523213357,1.8261047654783293),SE*labelscalefactor);
label("$r_1 - r_2$",(-0.648404598668290756,3.1317431971618523),SE*labelscalefactor);
/* dots and labels */
dot(O_1,dotstyle);
label("$O_1$", (0.0629524727693897,0.1494290064365599), NE * labelscalefactor);
dot(O_2,dotstyle);
label("$O_2$", (6.056676602422158,0.1494290064365599), NE * labelscalefactor);
//dot((1.3333333333333333,-3.7712361663282534), dotstyle);
dot(P_1, dotstyle);
//dot((4.,0.), dotstyle);
//dot((4.,0.), dotstyle);
dot(P_1, dotstyle);
label("$P_1$", (1.3868090159693944,3.8882326504189995), NE * labelscalefactor);
dot(P_2, dotstyle);
label("$P_2$", (6.72587881107271,1.9970090172761312), NE * labelscalefactor);
dot(Q, dotstyle);
label("$Q$", (0.7176068073188426,1.9970090172761312), NE * labelscalefactor);
clip((xmin,ymin)--(xmin,ymax)--(xmax,ymax)--(xmax,ymin)--cycle);
/* end of picture */
\end{asy}
\end{center}

We notice that since $O_1 P_1$ and $O_2 P_2$ are radii they meet the
common tangent at right angles, and in particular $O_1 P_1 \parallel
O_2 P_2$. Now it is clear to see that $QP_2 = r_1 + r_2$ and $QP_1 =
r_1 - r_2$. Now owing to Pythagoras,
$$
P_1P_2 = \sqrt{(r_1 + r_2)^2 - (r_1 - r_2)^2} = 2\sqrt{r_1 r_2}.
\ \square
$$

(b) Notice that due to the parallel conditions $O_1P_1P_2O_2$ is a trapezium,
with area
$$
\frac{r_1 + r_2}{2}\left(2\sqrt{r_1 r_2}\right).
$$
Now we just need to subtract off the two sectors. Let
$\angle P_1 O_1 O_2 = \theta$. Notice that $\angle O_1 O_2 P_2 = 180 - \theta$ due
to corresponding angles. Taking $\theta$ to be measured in radians gives the
area of the two sectors to be
$$
\frac12 r_1^2 \theta + \frac12 r_2^2 (\pi - \theta) =
\frac12 \theta \left(r_1^2 - r_2^2\right) + \frac12 \pi r_2^2.
$$
So the final area is
$$
(r_1 + r_2)\sqrt{r_1 r_2} - \frac12 \theta \left(r_1^2 - r_2^2\right) +
\frac12 \pi r_2^2. \ \blacksquare
$$

\textit{Comments: I first came across the fact (a) from a BPRP
(Black pen red pen) video, and then I actually used this fact to solve a problem
in a Year 12/13 geometry topic test. Part (b) actually came up as a question
with special case $r_1 = 3, r_2 = 1$ when
I was sitting the 2026 Jan TMUA, I was unable to solve it during contest.}
\end{solution}

\end{document}
